\section{Research and Authorship Ethics}\label{sec:research}

\subsection{Permitted research uses}
Researchers may use AI to explore literature, draft and edit text, write and debug analysis code, summarize and organize notes, and generate ideas, subject to journal, funder, and ethics-committee requirements, which prevail when stricter.

\subsection{Authorship and accountability}
\begin{enumerate}
  \item AI systems cannot be authors; human authors take full responsibility for content, accuracy, originality, and citations. \cmo{} requires clarity on human scholarly responsibility and accountability and allows institutions to specify whether and how AI tools may be acknowledged, provided acknowledgment does not displace human responsibility (Sec.~12). Publisher and funder policies apply in addition and should be checked for each target venue.
  \item Disclose substantive AI use (e.g., generation of text, analysis, or images) in the methods or acknowledgments, naming the tool and how it was used. Routine spelling or grammar assistance need not be disclosed unless the publisher requires it.
  \item Verify every AI-supplied reference and claim against the primary source; fabricated references are research misconduct when published.
  \item Maintain research records sufficient for reproducibility, including AI-assisted analysis steps.
\end{enumerate}

\subsection{Confidentiality in research and review}
Do not upload confidential manuscripts, grant proposals, peer-review materials, interview transcripts, or participant data to external AI tools without authorization and ethics clearance (Section~\ref{sec:privacy}).

\subsection{Student theses and dissertations}
Disclosure of AI use in theses and dissertations is required (\cmo{}, Sec.~11), using the format in Appendix~\ref{app:declarations}. Programs may set AI-use levels for thesis work; defaults follow Section~\ref{sec:integrity}. Advisers and students agree in writing at the outset on permitted AI use and disclosure, and the student remains responsible for authorship and defense of the work.

\subsection{Research ethics review}
Research that develops or deploys AI systems affecting human participants, or that uses AI on human-subject data, follows existing research ethics review, with the Tier~C pathway applying to high-risk systems (Section~\ref{sec:tools}). Authorship and intellectual property follow existing University policies.
